\usepackage{hyperref}

\usepackage{cleveref}
\usepackage{etoolbox}
\newtoggle{fullversion}
%set to long version
\toggletrue{fullversion}
% Set to proceeding version
%\togglefalse{fullversion}
\usepackage{adjustbox}
\pagenumbering{arabic}
\usepackage{caption}
\usepackage{scalerel}
\usepackage{graphicx}
\usepackage{xcolor}
\usepackage{tcolorbox}
\usepackage{amsmath,amsfonts}
\usepackage{ifsym}
\usepackage{mdframed}
\usepackage{listings}
\usepackage{wrapfig}
%\usepackage{amssymb}
%\usepackage{savetrees}

\renewcommand{\baselinestretch}{0.985}


\makeatletter
\let\c@lofdepth\relax
\let\c@lotdepth\relax
\makeatother
\usepackage{subfigure}
\tcbset{before={\par\pagebreak[0]\smallskip\parindent=0pt},after={\par\noindent}}
\usepackage{tree-dvips}
\usepackage{multirow}
%\newtheorem{mydef}{Definition}
\newtheorem{mytheorem}{Theorem}
\newtheorem{myexample}{Example}
\usepackage{qtree}
% issue-specific 
\newcommand{\sys}{\textsc{BLIP}\xspace}
\usepackage[ruled,linesnumbered,vlined]{algorithm2e}
\usepackage{algorithmic}
\usepackage{textcomp}
\usepackage{tikz}
\newcommand*\circled[1]{\tikz[baseline=(char.base)]{
    \node[shape=circle,black,draw,inner sep=0.5pt] (char) {#1};}}
\usetikzlibrary{positioning, shapes.geometric, arrows}
\usetikzlibrary{shapes, arrows.meta, fit, shapes.geometric, fit, backgrounds}
\definecolor{codegreen}{rgb}{0,0.6,0} % a darker green
\definecolor{darkblue}{rgb}{0.0, 0.0, 0.5}
\definecolor{darkblueilp}{RGB}{0, 0, 175}
\definecolor{darkred}{rgb}{0.5, 0.0, 0.0}
\definecolor{darkgreen}{rgb}{0.0, 0.5, 0.0}
\definecolor{lightgray}{gray}{0.95}
\newcommand{\appendixtext}[1]{#1}

\newcommand{\Hdash}{
  \rotatebox{90}{$\Vdash$}
}

\newcommand{\nHdash}{
  \rotatebox{90}{$\nVdash$}
}

\tcbuselibrary{listings, skins}
\usepackage[disable]{todonotes}
%\usepackage{minted}
\usepackage[T1]{fontenc}
\usepackage{listings}

\lstset{
  language=Python,           % Set the language for syntax highlighting
  basicstyle=\ttfamily\small, % Use monospace font with smaller size
  keywordstyle=\color{blue}, % Highlight keywords in blue
  stringstyle=\color{red},   % Highlight strings in red
  commentstyle=\color{gray}, % Highlight comments in gray
  breaklines=true,           % Automatically break long lines
  numbers=none,              % Disable line numbers
  frame=single,              % Add a single-line box around the code
  framerule=0.5pt            % Thickness of the frame (adjust as needed)
}
%
%\setminted{style=friendly, breaklines, autogobble, fontsize=\small}



\newcommand{\topic}[1]{\vspace{.5pt} \noindent{\bf #1.}}
\newcommand{\subtopic}[1]{\vspace{.5pt} \noindent{\em #1.}}

\newif\ifshowblock

\newcommand{\vldb}[1]{#1}
\newcommand{\sigmod}[1]{}

%uncomment for TR
\newcommand{\techreportsp}[1]{#1}
\newcommand{\techreport}[1]{#1}
\newcommand{\trs}[1]{#1\ignorespaces}
\newcommand{\papertext}[1]{}
% use {\color{..}#1} (not \textcolor) so these accept multi-paragraph arguments
\newcommand{\rone}[1]{{\color{black}#1}}
\newcommand{\rtwo}[1]{{\color{black}#1}}
\newcommand{\rthree}[1]{{\color{black}#1}}
\newcommand{\rmix}[1]{{\color{black}#1}}
\showblocktrue

%uncomment for paper
% \newcommand{\techreportsp}[1]{}
% \newcommand{\techreport}[1]{}
% \newcommand{\trs}[1]{\ignorespaces}
% \newcommand{\papertext}[1]{#1}
% \newcommand{\rone}[1]{\textcolor{black}{#1}}
% \newcommand{\rtwo}[1]{\textcolor{black}{#1}}
% \newcommand{\rthree}[1]{\textcolor{black}{#1}}
% \newcommand{\rmix}[1]{\textcolor{black}{#1}}
% \newcommand{\rone}[1]{\textcolor{blue}{#1}}
% \newcommand{\rtwo}[1]{\textcolor{myForestGreen}{#1}}
% \newcommand{\rthree}[1]{\textcolor{orange}{#1}}
% \newcommand{\rmix}[1]{\textcolor{magenta}{#1}}


\newcommand{\cover}[1]{}

\definecolor{myForestGreen}{RGB}{34,139,34}

%\usepackage{amssymb}
\usepackage[normalem]{ulem}
\newcommand{\delete}[1]{\sout{}}
\newcommand{\cradd}[1]{\textcolor{black}{#1}}



\newcommand{\agp}[1]{\textcolor{blue}{Aditya: #1}}
\newcommand{\compare}[1]{\textcolor{red}{#1}}
\newcommand{\sep}[1]{\textcolor{red}{Sep: #1}}
\newcommand{\update}[1]{\textcolor{black}{#1}}
\newcommand{\yiming}[1]{\textcolor{orange}{Yiming: #1}}
\newcommand{\edit}[1]{\textcolor{black}{#1}}
%\newcommand{\anonymous}[2]{#1}% anon version
 \newcommand{\anonymous}[2]{#2} % non-anon version
\newcommand{\revised}[1]{\textcolor{black}{#1}}
\newcommand{\newexp}[1]{\textcolor{blue}{#1}}
\newcommand{\question}[1]{\textcolor{gray}{#1}}

\usepackage{xcolor}
\usepackage{tcolorbox}

\usepackage[normalem]{ulem}
\newcommand{\stext}[1]{\sout{#1}}


\usepackage{tikz}

\newcommand{\mstrike}[1]{%
\tikz[baseline=(X.base)]{
  \node[inner sep=0pt, outer sep=0pt] (X) {$\displaystyle #1$};
  \draw[line width=0.5pt] (X.west) -- (X.east);
}%
}

\newcommand{\response}[1]{\noindent\textbf{\textgreater\ #1}}
\newcounter{definition}
\newenvironment{definition}[1][]{\refstepcounter{definition}\par\smallskip\textsc{Definition~\thedefinition.\ #1}}{\smallskip}

\newcounter{problem}
\newenvironment{problem}[1][]{\refstepcounter{problem}\par\smallskip\textsc{Problem~\theproblem.\ #1}}{\smallskip}


\newcounter{assumption}
\newenvironment{assumption}[1][]{\refstepcounter{assumption}\par\smallskip\textsc{Assumption~\theassumption.\ #1}}{\smallskip}

\newcommand{\code}[1]{\texttt{\small #1}}

\newcounter{example}
\newenvironment{example}[1][]{\refstepcounter{example}\par\smallskip\textsc{Example~\theexample.\ #1}}{$\square$\smallskip}

\newcounter{theorem}
\newenvironment{theorem}[1][]{\refstepcounter{theorem}\par\smallskip\textsc{Theorem~\thetheorem.\ #1}}{\smallskip}

\newcounter{proposition}
\newenvironment{proposition}[1][]{
    \refstepcounter{proposition}
    \par\smallskip
    \textsc{Proposition~\theproposition.} \textit{#1}%\par\smallskip
}{
    \smallskip
}


\usepackage{enumitem}
\newcommand{\squishlist}{
	\begin{list}{$\bullet$}
		{
			\setlength{\itemsep}{0pt}
			\setlength{\parsep}{1pt}
			\setlength{\topsep}{1pt}
			\setlength{\partopsep}{0pt}
			\setlength{\leftmargin}{1.5em}
			\setlength{\labelwidth}{1em}
			\setlength{\labelsep}{0.5em} } }
	
\newcommand{\squishend}{\end{list}}

\newcommand{\squishlistnum}{
	\begin{enumerate}
		{
			\setlength{\itemsep}{0pt}
			\setlength{\parsep}{1pt}
			\setlength{\topsep}{1pt}
			\setlength{\partopsep}{0pt}
			\setlength{\leftmargin}{1.5em}
			\setlength{\labelwidth}{1em}
			\setlength{\labelsep}{0.5em}
		}
}

\newcommand{\squishendnum}{\end{enumerate}}
\usepackage{multirow}

\usepackage{amsthm}
%\newtheorem{assumption}{Assumption}
\usepackage{mathrsfs}

% \usepackage{tikz}
% \newcommand*\circled[1]{\tikz[baseline=(char.base)]{
%     \node[shape=circle,red,draw,inner sep=0.5pt] (char) {#1};}}


\newcommand\vldbdoi{10.14778/3836663.3836668}
\newcommand\vldbpages{2992 - 3005}
% issue-specific
\newcommand\vldbvolume{19}
\newcommand\vldbissue{11}
\newcommand\vldbyear{2026}
% should be fine as it is
\newcommand\vldbauthors{\authors}
\newcommand\vldbtitle{\shorttitle} 
% leave empty if no availability url should be set
\newcommand\vldbavailabilityurl{https://github.com/yiminl18/BLIP.git}
% whether page numbers should be shown or not, use 'plain' for review versions, 'empty' for camera ready
\newcommand\vldbpagestyle{empty}


\lstdefinestyle{SQLStyle}{
  language=SQL,
  showspaces=false,
  basicstyle=\ttfamily\footnotesize,
  commentstyle=\color{gray},
  mathescape=true,
  numbers=none,
  escapeinside={^}{^},
  captionpos=b,
  mathescape=false,
}
\newcommand{\mypython}[1]{%
  \begin{tcolorbox}[
    colback=gray!10,
    colframe=black,
    boxrule=0.5pt,
    arc=2pt,
    left=5pt,
    right=5pt,
    top=2pt,
    bottom=2pt
  ]
  \small\ttfamily #1
  \end{tcolorbox}
}

\NewDocumentCommand{\sidequote}{m +m O{black}}{%
 \begin{tcolorbox}[
enhanced,
breakable,
frame hidden,
interior hidden,
sharp corners,
borderline west={2.5pt}{0pt}{gray!50},
left=15pt,
right=0pt,
top=0pt,
bottom=0pt,
coltext=#3 % <--- Applies the color safely across page boundaries
]
{\normalcolor\small\sffamily\bfseries #2}\par\medskip % The Source
#1 % The Quoted Text
\end{tcolorbox}
}
 